En una exposició d'art hi ha una obra que consisteix en una moneda situada a l'interior d'un bloc massís de policarbonat transparent, just al centre de la cara posterior, tal com indica la figura. El bloc té forma d'ortoedre (de capsa de sabates) i un índex de refracció d'1,58. Una persona està observant l'obra des del punt mitjà de la cara oposada i arriba una segona persona i se situa a la dreta de la primera. Sorprenentment, no veu la moneda de l'interior del bloc.

\figura[0.45]{fig1.png}

\begin{apartat}{1,25}
A partir de la llei de Snell, deduïu l'expressió de l'angle límit (o angle crític) en funció dels índexs de refracció dels dos medis. Calculeu l'angle límit amb les dades del problema. Justifiqueu si es podria donar aquest fenomen en el cas que s'invertissin els medis.
\end{apartat}

\begin{apartat}{1,25}
A quina distància màxima, $d$, s'hauria de col·locar la segona persona respecte de la primera per veure la moneda? Considereu un raig de llum que surt de la moneda i arriba a la interfície de policarbonat-aire amb aquest angle límit i dibuixeu un esquema dels raigs incident, reflectit i refractat per a aquest cas.
\end{apartat}

\dades{Índex de refracció de l'aire $= 1$.}
